\subsection{Hasil dan Perbandingan dengan \textit{Scikit-learn}}
\label{subsec:dtl-hasil}

\todo{Bandingkan dengan \texttt{DecisionTreeClassifier} dalam tiga konfigurasi wajib dari spesifikasi: ID3 (\texttt{criterion='entropy'}), C4.5 (\texttt{criterion='entropy'} + \texttt{ccp\_alpha}), dan CART (\texttt{criterion='gini'}). Jelaskan \textit{preprocessing} tambahan untuk \textit{sklearn} (\textit{encoding} dan imputasi) serta cara memilih \texttt{ccp\_alpha}. Pakai \textit{preprocessing} dan \textit{fold} yang sama agar perbedaan hanya berasal dari implementasi \cite{pedregosa2011sklearn}.}

\begin{table}[H]
	\centering
	\caption{Perbandingan DTL \textit{from scratch} dan \textit{scikit-learn}}
	\label{tab:dtl-hasil}
	\begin{tabular}{l c c c c c}
		\toprule
		\textbf{Implementasi} & \textbf{Macro F1} & \textbf{Accuracy} & \textbf{Precision} & \textbf{Recall} & \textbf{Waktu (s)} \\
		\midrule
		\textit{From scratch} (C4.5) & & & & & \\
		\textit{Sklearn} ID3 & & & & & \\
		\textit{Sklearn} C4.5 (\texttt{ccp\_alpha}) & & & & & \\
		\textit{Sklearn} CART & & & & & \\
		\bottomrule
	\end{tabular}
\end{table}

\begin{figure}[H]
	\centering
	\placeholdergambar{\textit{Confusion matrix} \textit{from scratch} vs \textit{scikit-learn}}
	\caption{\textit{Confusion matrix} DTL}
	\label{fig:dtl-cm}
\end{figure}

\subsubsection{Pengaruh \textit{Pruning}}
\todo{Bandingkan strategi \textit{pruning} (tanpa, \textit{pessimistic}, \textit{reduced error}, \textit{cost-complexity}): kedalaman, jumlah daun, dan \textit{macro F1}. Sertakan grafik.}

\subsubsection{Analisis}
\todo{Bandingkan kedalaman, jumlah daun, dan kinerja. Jelaskan perbedaan karena \textit{gain ratio} vs entropi/gini, \textit{multiway} vs \textit{binary split}, penanganan \textit{null}, dan jenis \textit{pruning}.}
